% ============================================================
% Pain/Gain-Share Vergütungsmodell — TikZ-Grafik
% ============================================================
% EINBINDUNG: % ============================================================
% Pain/Gain-Share Vergütungsmodell — TikZ-Grafik
% ============================================================
% EINBINDUNG: % ============================================================
% Pain/Gain-Share Vergütungsmodell — TikZ-Grafik
% ============================================================
% EINBINDUNG: % ============================================================
% Pain/Gain-Share Vergütungsmodell — TikZ-Grafik
% ============================================================
% EINBINDUNG: \input{05_figures/Pain_Gain_Share.tex}
% Nutzt ausschließlich tu-green aus settings.tex + xcolor-Shades.
% ============================================================

\begin{figure}[htbp]
    \centering
    \fcolorbox{gray!50}{white}{%
        \begin{minipage}{\dimexpr\textwidth-2\fboxsep-2\fboxrule\relax}
            \centering
            \vspace{0.5cm}
            \begin{tikzpicture}[
                every node/.style={font=\sffamily},
                seglabel/.style={font=\sffamily\itshape\scriptsize, align=center},
                collabel/.style={font=\sffamily\footnotesize\bfseries, align=center}
            ]

            % ===== Konfiguration =====
            \def\barW{2.8}       % Balkenbreite
            \def\gap{1.0}        % Lücke zwischen Balken

            % X-Positionen (linke Kante jedes Balkens)
            \def\xA{0}
            \pgfmathsetmacro{\xB}{\barW + \gap}
            \pgfmathsetmacro{\xC}{2*(\barW + \gap)}
            \pgfmathsetmacro{\xD}{3*(\barW + \gap)}

            % ===== Höhen (in cm) =====
            % Fall 1: DK=5.0  GK=1.2  G=0.8  Total=7.0
            % Fall 2: DK=5.8  GK=0.7  G=0.5  Total=7.0
            % Fall 3: DK=3.6  GK=1.2  G=0.8  B1=0.7  B2=0.7  Total=7.0
            % Fall 4: DK=8.5

            % ============================================
            % FALL 1: Basis-Zielkosten
            % ============================================
            \fill[tu-green]       (\xA, 0)   rectangle +(\barW, 5.0);
            \fill[tu-green!65]    (\xA, 5.0) rectangle +(\barW, 1.2);
            \fill[tu-green!50]    (\xA, 6.2) rectangle +(\barW, 0.8);

            \node[seglabel, white]   at (\xA+\barW/2, 2.5)  {Direkte\\Kosten};
            \node[seglabel, white]   at (\xA+\barW/2, 5.6)  {GK-Zuschlag};
            \node[seglabel]          at (\xA+\barW/2, 6.6)  {Gewinn};

            \node[collabel] at (\xA+\barW/2, -0.7) {Fall 1:\\Basis-Zielkosten};

            % ============================================
            % FALL 2: Mehrkosten (Pain)
            % ============================================
            \fill[tu-green]       (\xB, 0)   rectangle +(\barW, 5.8);
            \fill[tu-green!65]    (\xB, 5.8) rectangle +(\barW, 0.7);
            \fill[tu-green!50]    (\xB, 6.5) rectangle +(\barW, 0.5);

            \node[seglabel, white]   at (\xB+\barW/2, 2.9)  {Direkte\\Kosten};
            \node[seglabel, white]   at (\xB+\barW/2, 6.15) {GK-Zuschlag};
            \node[seglabel]          at (\xB+\barW/2, 6.75) {Gewinn};

            \node[collabel] at (\xB+\barW/2, -0.7) {Fall 2:\\Mehrkosten};

            % ============================================
            % FALL 3: Kosteneinsparung (Gain)
            % ============================================
            \fill[tu-green]       (\xC, 0)   rectangle +(\barW, 3.6);
            \fill[tu-green!65]    (\xC, 3.6) rectangle +(\barW, 1.2);
            \fill[tu-green!50]    (\xC, 4.8) rectangle +(\barW, 0.8);
            \fill[tu-green!25]    (\xC, 5.6) rectangle +(\barW, 0.7);
            \fill[tu-green!25]    (\xC, 6.3) rectangle +(\barW, 0.7);

            \node[seglabel, white]   at (\xC+\barW/2, 1.8)  {Direkte\\Kosten};
            \node[seglabel, white]   at (\xC+\barW/2, 4.2)  {GK-Zuschlag};
            \node[seglabel]          at (\xC+\barW/2, 5.2)  {Gewinn};
            \node[seglabel]          at (\xC+\barW/2, 5.95) {Bonusteilung};
            \node[seglabel]          at (\xC+\barW/2, 6.65) {Bonusteilung};

            \node[collabel] at (\xC+\barW/2, -0.7) {Fall 3:\\Kosteneinsparung};

            % ============================================
            % FALL 4: Tatsächliche Kosten
            % ============================================
            \fill[tu-green]       (\xD, 0)   rectangle +(\barW, 8.5);

            \node[seglabel, white]   at (\xD+\barW/2, 4.25) {Direkte\\Kosten};

            \node[collabel] at (\xD+\barW/2, -0.7) {Fall 4:\\Tatsächliche Kosten};

            % ============================================
            % LINIE 1: Oberkante Direkte Kosten (durchgehend)
            % ============================================
            \draw[black!65, dashed, line width=0.7pt]
                (\xA, 5.0) -- (\xA+\barW, 5.0) --
                (\xB, 5.8) -- (\xB+\barW, 5.8) --
                (\xC, 3.6) -- (\xC+\barW, 3.6) --
                (\xD, 8.5) -- (\xD+\barW, 8.5);

            % ============================================
            % LINIE 2: Oberkante Gewinn (F1–F3) → DK top F4
            % ============================================
            \draw[black!65, dashed, line width=0.7pt]
                (\xA, 7.0) -- (\xA+\barW, 7.0) --
                (\xB, 7.0) -- (\xB+\barW, 7.0) --
                (\xC, 5.6) -- (\xC+\barW, 5.6) --
                (\xD, 8.5) -- (\xD+\barW, 8.5);

            \end{tikzpicture}
            \vspace{0.4cm}
        \end{minipage}%
    }
    \caption[Vergütungsstufen in der IPA]{Vergütungsstufen und Pain/Gain-Share-Mechanismus in der IPA (Quelle: Eigene Darstellung in Anlehnung an \textcite{becker_integrierte_2022}).}
    \label{fig:pain_gain_share}
\end{figure}

% Nutzt ausschließlich tu-green aus settings.tex + xcolor-Shades.
% ============================================================

\begin{figure}[htbp]
    \centering
    \fcolorbox{gray!50}{white}{%
        \begin{minipage}{\dimexpr\textwidth-2\fboxsep-2\fboxrule\relax}
            \centering
            \vspace{0.5cm}
            \begin{tikzpicture}[
                every node/.style={font=\sffamily},
                seglabel/.style={font=\sffamily\itshape\scriptsize, align=center},
                collabel/.style={font=\sffamily\footnotesize\bfseries, align=center}
            ]

            % ===== Konfiguration =====
            \def\barW{2.8}       % Balkenbreite
            \def\gap{1.0}        % Lücke zwischen Balken

            % X-Positionen (linke Kante jedes Balkens)
            \def\xA{0}
            \pgfmathsetmacro{\xB}{\barW + \gap}
            \pgfmathsetmacro{\xC}{2*(\barW + \gap)}
            \pgfmathsetmacro{\xD}{3*(\barW + \gap)}

            % ===== Höhen (in cm) =====
            % Fall 1: DK=5.0  GK=1.2  G=0.8  Total=7.0
            % Fall 2: DK=5.8  GK=0.7  G=0.5  Total=7.0
            % Fall 3: DK=3.6  GK=1.2  G=0.8  B1=0.7  B2=0.7  Total=7.0
            % Fall 4: DK=8.5

            % ============================================
            % FALL 1: Basis-Zielkosten
            % ============================================
            \fill[tu-green]       (\xA, 0)   rectangle +(\barW, 5.0);
            \fill[tu-green!65]    (\xA, 5.0) rectangle +(\barW, 1.2);
            \fill[tu-green!50]    (\xA, 6.2) rectangle +(\barW, 0.8);

            \node[seglabel, white]   at (\xA+\barW/2, 2.5)  {Direkte\\Kosten};
            \node[seglabel, white]   at (\xA+\barW/2, 5.6)  {GK-Zuschlag};
            \node[seglabel]          at (\xA+\barW/2, 6.6)  {Gewinn};

            \node[collabel] at (\xA+\barW/2, -0.7) {Fall 1:\\Basis-Zielkosten};

            % ============================================
            % FALL 2: Mehrkosten (Pain)
            % ============================================
            \fill[tu-green]       (\xB, 0)   rectangle +(\barW, 5.8);
            \fill[tu-green!65]    (\xB, 5.8) rectangle +(\barW, 0.7);
            \fill[tu-green!50]    (\xB, 6.5) rectangle +(\barW, 0.5);

            \node[seglabel, white]   at (\xB+\barW/2, 2.9)  {Direkte\\Kosten};
            \node[seglabel, white]   at (\xB+\barW/2, 6.15) {GK-Zuschlag};
            \node[seglabel]          at (\xB+\barW/2, 6.75) {Gewinn};

            \node[collabel] at (\xB+\barW/2, -0.7) {Fall 2:\\Mehrkosten};

            % ============================================
            % FALL 3: Kosteneinsparung (Gain)
            % ============================================
            \fill[tu-green]       (\xC, 0)   rectangle +(\barW, 3.6);
            \fill[tu-green!65]    (\xC, 3.6) rectangle +(\barW, 1.2);
            \fill[tu-green!50]    (\xC, 4.8) rectangle +(\barW, 0.8);
            \fill[tu-green!25]    (\xC, 5.6) rectangle +(\barW, 0.7);
            \fill[tu-green!25]    (\xC, 6.3) rectangle +(\barW, 0.7);

            \node[seglabel, white]   at (\xC+\barW/2, 1.8)  {Direkte\\Kosten};
            \node[seglabel, white]   at (\xC+\barW/2, 4.2)  {GK-Zuschlag};
            \node[seglabel]          at (\xC+\barW/2, 5.2)  {Gewinn};
            \node[seglabel]          at (\xC+\barW/2, 5.95) {Bonusteilung};
            \node[seglabel]          at (\xC+\barW/2, 6.65) {Bonusteilung};

            \node[collabel] at (\xC+\barW/2, -0.7) {Fall 3:\\Kosteneinsparung};

            % ============================================
            % FALL 4: Tatsächliche Kosten
            % ============================================
            \fill[tu-green]       (\xD, 0)   rectangle +(\barW, 8.5);

            \node[seglabel, white]   at (\xD+\barW/2, 4.25) {Direkte\\Kosten};

            \node[collabel] at (\xD+\barW/2, -0.7) {Fall 4:\\Tatsächliche Kosten};

            % ============================================
            % LINIE 1: Oberkante Direkte Kosten (durchgehend)
            % ============================================
            \draw[black!65, dashed, line width=0.7pt]
                (\xA, 5.0) -- (\xA+\barW, 5.0) --
                (\xB, 5.8) -- (\xB+\barW, 5.8) --
                (\xC, 3.6) -- (\xC+\barW, 3.6) --
                (\xD, 8.5) -- (\xD+\barW, 8.5);

            % ============================================
            % LINIE 2: Oberkante Gewinn (F1–F3) → DK top F4
            % ============================================
            \draw[black!65, dashed, line width=0.7pt]
                (\xA, 7.0) -- (\xA+\barW, 7.0) --
                (\xB, 7.0) -- (\xB+\barW, 7.0) --
                (\xC, 5.6) -- (\xC+\barW, 5.6) --
                (\xD, 8.5) -- (\xD+\barW, 8.5);

            \end{tikzpicture}
            \vspace{0.4cm}
        \end{minipage}%
    }
    \caption[Vergütungsstufen in der IPA]{Vergütungsstufen und Pain/Gain-Share-Mechanismus in der IPA (Quelle: Eigene Darstellung in Anlehnung an \textcite{becker_integrierte_2022}).}
    \label{fig:pain_gain_share}
\end{figure}

% Nutzt ausschließlich tu-green aus settings.tex + xcolor-Shades.
% ============================================================

\begin{figure}[htbp]
    \centering
    \fcolorbox{gray!50}{white}{%
        \begin{minipage}{\dimexpr\textwidth-2\fboxsep-2\fboxrule\relax}
            \centering
            \vspace{0.5cm}
            \begin{tikzpicture}[
                every node/.style={font=\sffamily},
                seglabel/.style={font=\sffamily\itshape\scriptsize, align=center},
                collabel/.style={font=\sffamily\footnotesize\bfseries, align=center}
            ]

            % ===== Konfiguration =====
            \def\barW{2.8}       % Balkenbreite
            \def\gap{1.0}        % Lücke zwischen Balken

            % X-Positionen (linke Kante jedes Balkens)
            \def\xA{0}
            \pgfmathsetmacro{\xB}{\barW + \gap}
            \pgfmathsetmacro{\xC}{2*(\barW + \gap)}
            \pgfmathsetmacro{\xD}{3*(\barW + \gap)}

            % ===== Höhen (in cm) =====
            % Fall 1: DK=5.0  GK=1.2  G=0.8  Total=7.0
            % Fall 2: DK=5.8  GK=0.7  G=0.5  Total=7.0
            % Fall 3: DK=3.6  GK=1.2  G=0.8  B1=0.7  B2=0.7  Total=7.0
            % Fall 4: DK=8.5

            % ============================================
            % FALL 1: Basis-Zielkosten
            % ============================================
            \fill[tu-green]       (\xA, 0)   rectangle +(\barW, 5.0);
            \fill[tu-green!65]    (\xA, 5.0) rectangle +(\barW, 1.2);
            \fill[tu-green!50]    (\xA, 6.2) rectangle +(\barW, 0.8);

            \node[seglabel, white]   at (\xA+\barW/2, 2.5)  {Direkte\\Kosten};
            \node[seglabel, white]   at (\xA+\barW/2, 5.6)  {GK-Zuschlag};
            \node[seglabel]          at (\xA+\barW/2, 6.6)  {Gewinn};

            \node[collabel] at (\xA+\barW/2, -0.7) {Fall 1:\\Basis-Zielkosten};

            % ============================================
            % FALL 2: Mehrkosten (Pain)
            % ============================================
            \fill[tu-green]       (\xB, 0)   rectangle +(\barW, 5.8);
            \fill[tu-green!65]    (\xB, 5.8) rectangle +(\barW, 0.7);
            \fill[tu-green!50]    (\xB, 6.5) rectangle +(\barW, 0.5);

            \node[seglabel, white]   at (\xB+\barW/2, 2.9)  {Direkte\\Kosten};
            \node[seglabel, white]   at (\xB+\barW/2, 6.15) {GK-Zuschlag};
            \node[seglabel]          at (\xB+\barW/2, 6.75) {Gewinn};

            \node[collabel] at (\xB+\barW/2, -0.7) {Fall 2:\\Mehrkosten};

            % ============================================
            % FALL 3: Kosteneinsparung (Gain)
            % ============================================
            \fill[tu-green]       (\xC, 0)   rectangle +(\barW, 3.6);
            \fill[tu-green!65]    (\xC, 3.6) rectangle +(\barW, 1.2);
            \fill[tu-green!50]    (\xC, 4.8) rectangle +(\barW, 0.8);
            \fill[tu-green!25]    (\xC, 5.6) rectangle +(\barW, 0.7);
            \fill[tu-green!25]    (\xC, 6.3) rectangle +(\barW, 0.7);

            \node[seglabel, white]   at (\xC+\barW/2, 1.8)  {Direkte\\Kosten};
            \node[seglabel, white]   at (\xC+\barW/2, 4.2)  {GK-Zuschlag};
            \node[seglabel]          at (\xC+\barW/2, 5.2)  {Gewinn};
            \node[seglabel]          at (\xC+\barW/2, 5.95) {Bonusteilung};
            \node[seglabel]          at (\xC+\barW/2, 6.65) {Bonusteilung};

            \node[collabel] at (\xC+\barW/2, -0.7) {Fall 3:\\Kosteneinsparung};

            % ============================================
            % FALL 4: Tatsächliche Kosten
            % ============================================
            \fill[tu-green]       (\xD, 0)   rectangle +(\barW, 8.5);

            \node[seglabel, white]   at (\xD+\barW/2, 4.25) {Direkte\\Kosten};

            \node[collabel] at (\xD+\barW/2, -0.7) {Fall 4:\\Tatsächliche Kosten};

            % ============================================
            % LINIE 1: Oberkante Direkte Kosten (durchgehend)
            % ============================================
            \draw[black!65, dashed, line width=0.7pt]
                (\xA, 5.0) -- (\xA+\barW, 5.0) --
                (\xB, 5.8) -- (\xB+\barW, 5.8) --
                (\xC, 3.6) -- (\xC+\barW, 3.6) --
                (\xD, 8.5) -- (\xD+\barW, 8.5);

            % ============================================
            % LINIE 2: Oberkante Gewinn (F1–F3) → DK top F4
            % ============================================
            \draw[black!65, dashed, line width=0.7pt]
                (\xA, 7.0) -- (\xA+\barW, 7.0) --
                (\xB, 7.0) -- (\xB+\barW, 7.0) --
                (\xC, 5.6) -- (\xC+\barW, 5.6) --
                (\xD, 8.5) -- (\xD+\barW, 8.5);

            \end{tikzpicture}
            \vspace{0.4cm}
        \end{minipage}%
    }
    \caption[Vergütungsstufen in der IPA]{Vergütungsstufen und Pain/Gain-Share-Mechanismus in der IPA (Quelle: Eigene Darstellung in Anlehnung an \textcite{becker_integrierte_2022}).}
    \label{fig:pain_gain_share}
\end{figure}

% Nutzt ausschließlich tu-green aus settings.tex + xcolor-Shades.
% ============================================================

\begin{figure}[htbp]
    \centering
    \fcolorbox{gray!50}{white}{%
        \begin{minipage}{\dimexpr\textwidth-2\fboxsep-2\fboxrule\relax}
            \centering
            \vspace{0.5cm}
            \begin{tikzpicture}[
                every node/.style={font=\sffamily},
                seglabel/.style={font=\sffamily\itshape\scriptsize, align=center},
                collabel/.style={font=\sffamily\footnotesize\bfseries, align=center}
            ]

            % ===== Konfiguration =====
            \def\barW{2.8}       % Balkenbreite
            \def\gap{1.0}        % Lücke zwischen Balken

            % X-Positionen (linke Kante jedes Balkens)
            \def\xA{0}
            \pgfmathsetmacro{\xB}{\barW + \gap}
            \pgfmathsetmacro{\xC}{2*(\barW + \gap)}
            \pgfmathsetmacro{\xD}{3*(\barW + \gap)}

            % ===== Höhen (in cm) =====
            % Fall 1: DK=5.0  GK=1.2  G=0.8  Total=7.0
            % Fall 2: DK=5.8  GK=0.7  G=0.5  Total=7.0
            % Fall 3: DK=3.6  GK=1.2  G=0.8  B1=0.7  B2=0.7  Total=7.0
            % Fall 4: DK=8.5

            % ============================================
            % FALL 1: Basis-Zielkosten
            % ============================================
            \fill[tu-green]       (\xA, 0)   rectangle +(\barW, 5.0);
            \fill[tu-green!65]    (\xA, 5.0) rectangle +(\barW, 1.2);
            \fill[tu-green!50]    (\xA, 6.2) rectangle +(\barW, 0.8);

            \node[seglabel, white]   at (\xA+\barW/2, 2.5)  {Direkte\\Kosten};
            \node[seglabel, white]   at (\xA+\barW/2, 5.6)  {GK-Zuschlag};
            \node[seglabel]          at (\xA+\barW/2, 6.6)  {Gewinn};

            \node[collabel] at (\xA+\barW/2, -0.7) {Fall 1:\\Basis-Zielkosten};

            % ============================================
            % FALL 2: Mehrkosten (Pain)
            % ============================================
            \fill[tu-green]       (\xB, 0)   rectangle +(\barW, 5.8);
            \fill[tu-green!65]    (\xB, 5.8) rectangle +(\barW, 0.7);
            \fill[tu-green!50]    (\xB, 6.5) rectangle +(\barW, 0.5);

            \node[seglabel, white]   at (\xB+\barW/2, 2.9)  {Direkte\\Kosten};
            \node[seglabel, white]   at (\xB+\barW/2, 6.15) {GK-Zuschlag};
            \node[seglabel]          at (\xB+\barW/2, 6.75) {Gewinn};

            \node[collabel] at (\xB+\barW/2, -0.7) {Fall 2:\\Mehrkosten};

            % ============================================
            % FALL 3: Kosteneinsparung (Gain)
            % ============================================
            \fill[tu-green]       (\xC, 0)   rectangle +(\barW, 3.6);
            \fill[tu-green!65]    (\xC, 3.6) rectangle +(\barW, 1.2);
            \fill[tu-green!50]    (\xC, 4.8) rectangle +(\barW, 0.8);
            \fill[tu-green!25]    (\xC, 5.6) rectangle +(\barW, 0.7);
            \fill[tu-green!25]    (\xC, 6.3) rectangle +(\barW, 0.7);

            \node[seglabel, white]   at (\xC+\barW/2, 1.8)  {Direkte\\Kosten};
            \node[seglabel, white]   at (\xC+\barW/2, 4.2)  {GK-Zuschlag};
            \node[seglabel]          at (\xC+\barW/2, 5.2)  {Gewinn};
            \node[seglabel]          at (\xC+\barW/2, 5.95) {Bonusteilung};
            \node[seglabel]          at (\xC+\barW/2, 6.65) {Bonusteilung};

            \node[collabel] at (\xC+\barW/2, -0.7) {Fall 3:\\Kosteneinsparung};

            % ============================================
            % FALL 4: Tatsächliche Kosten
            % ============================================
            \fill[tu-green]       (\xD, 0)   rectangle +(\barW, 8.5);

            \node[seglabel, white]   at (\xD+\barW/2, 4.25) {Direkte\\Kosten};

            \node[collabel] at (\xD+\barW/2, -0.7) {Fall 4:\\Tatsächliche Kosten};

            % ============================================
            % LINIE 1: Oberkante Direkte Kosten (durchgehend)
            % ============================================
            \draw[black!65, dashed, line width=0.7pt]
                (\xA, 5.0) -- (\xA+\barW, 5.0) --
                (\xB, 5.8) -- (\xB+\barW, 5.8) --
                (\xC, 3.6) -- (\xC+\barW, 3.6) --
                (\xD, 8.5) -- (\xD+\barW, 8.5);

            % ============================================
            % LINIE 2: Oberkante Gewinn (F1–F3) → DK top F4
            % ============================================
            \draw[black!65, dashed, line width=0.7pt]
                (\xA, 7.0) -- (\xA+\barW, 7.0) --
                (\xB, 7.0) -- (\xB+\barW, 7.0) --
                (\xC, 5.6) -- (\xC+\barW, 5.6) --
                (\xD, 8.5) -- (\xD+\barW, 8.5);

            \end{tikzpicture}
            \vspace{0.4cm}
        \end{minipage}%
    }
    \caption[Vergütungsstufen in der IPA]{Vergütungsstufen und Pain/Gain-Share-Mechanismus in der IPA (Quelle: Eigene Darstellung in Anlehnung an \textcite{becker_integrierte_2022}).}
    \label{fig:pain_gain_share}
\end{figure}
